\section{La brecha entre las matemáticas de competencia y las matemáticas de frontera}

\subsection{Limitaciones actuales de la capacidad matemática de los LLM}

\begin{itemize}
    \item O1/O3 tienen un desempeño sobresaliente en competencias preparatorias como AMC/AIME
    \item El benchmark Frontier Math contiene problemas de matemáticas de nivel de investigación; O3 resolvió más del 20\% (pero los problemas están diseñados para tener respuestas numéricas)
    \item Evaluación de Terence Tao: O1 todavía tiene dificultades en tareas de investigación matemática de frontera
\end{itemize}

\subsection{El problema de la generación de demostraciones}

\begin{warningbox}{El grave problema de los LLM al escribir demostraciones}
\begin{itemize}
    \item Competencia Putnam (nivel licenciatura): ChatGPT obtiene en su mayoría solo 1--2/10 puntos
    \item Evaluación USAMO: todos los modelos obtienen una puntuación de alrededor del 5\%
    \item Tipo de error: errores «simples» como la inversión del sentido de una desigualdad, ocultos dentro de demostraciones largas y sumamente difíciles de detectar
\end{itemize}
Causa raíz: en las matemáticas de frontera los datos son \textbf{escasos} y las demostraciones son \textbf{no verificables}.
\end{warningbox}

